% Options for packages loaded elsewhere
\PassOptionsToPackage{unicode}{hyperref}
\PassOptionsToPackage{hyphens}{url}
\documentclass[
]{article}
\usepackage{xcolor}
\usepackage[margin=1in]{geometry}
\usepackage{amsmath,amssymb}
\setcounter{secnumdepth}{-\maxdimen} % remove section numbering
\usepackage{iftex}
\ifPDFTeX
  \usepackage[T1]{fontenc}
  \usepackage[utf8]{inputenc}
  \usepackage{textcomp} % provide euro and other symbols
\else % if luatex or xetex
  \usepackage{unicode-math} % this also loads fontspec
  \defaultfontfeatures{Scale=MatchLowercase}
  \defaultfontfeatures[\rmfamily]{Ligatures=TeX,Scale=1}
\fi
\usepackage{lmodern}
\ifPDFTeX\else
  % xetex/luatex font selection
\fi
% Use upquote if available, for straight quotes in verbatim environments
\IfFileExists{upquote.sty}{\usepackage{upquote}}{}
\IfFileExists{microtype.sty}{% use microtype if available
  \usepackage[]{microtype}
  \UseMicrotypeSet[protrusion]{basicmath} % disable protrusion for tt fonts
}{}
\makeatletter
\@ifundefined{KOMAClassName}{% if non-KOMA class
  \IfFileExists{parskip.sty}{%
    \usepackage{parskip}
  }{% else
    \setlength{\parindent}{0pt}
    \setlength{\parskip}{6pt plus 2pt minus 1pt}}
}{% if KOMA class
  \KOMAoptions{parskip=half}}
\makeatother
\usepackage{color}
\usepackage{fancyvrb}
\newcommand{\VerbBar}{|}
\newcommand{\VERB}{\Verb[commandchars=\\\{\}]}
\DefineVerbatimEnvironment{Highlighting}{Verbatim}{commandchars=\\\{\}}
% Add ',fontsize=\small' for more characters per line
\usepackage{framed}
\definecolor{shadecolor}{RGB}{248,248,248}
\newenvironment{Shaded}{\begin{snugshade}}{\end{snugshade}}
\newcommand{\AlertTok}[1]{\textcolor[rgb]{0.94,0.16,0.16}{#1}}
\newcommand{\AnnotationTok}[1]{\textcolor[rgb]{0.56,0.35,0.01}{\textbf{\textit{#1}}}}
\newcommand{\AttributeTok}[1]{\textcolor[rgb]{0.13,0.29,0.53}{#1}}
\newcommand{\BaseNTok}[1]{\textcolor[rgb]{0.00,0.00,0.81}{#1}}
\newcommand{\BuiltInTok}[1]{#1}
\newcommand{\CharTok}[1]{\textcolor[rgb]{0.31,0.60,0.02}{#1}}
\newcommand{\CommentTok}[1]{\textcolor[rgb]{0.56,0.35,0.01}{\textit{#1}}}
\newcommand{\CommentVarTok}[1]{\textcolor[rgb]{0.56,0.35,0.01}{\textbf{\textit{#1}}}}
\newcommand{\ConstantTok}[1]{\textcolor[rgb]{0.56,0.35,0.01}{#1}}
\newcommand{\ControlFlowTok}[1]{\textcolor[rgb]{0.13,0.29,0.53}{\textbf{#1}}}
\newcommand{\DataTypeTok}[1]{\textcolor[rgb]{0.13,0.29,0.53}{#1}}
\newcommand{\DecValTok}[1]{\textcolor[rgb]{0.00,0.00,0.81}{#1}}
\newcommand{\DocumentationTok}[1]{\textcolor[rgb]{0.56,0.35,0.01}{\textbf{\textit{#1}}}}
\newcommand{\ErrorTok}[1]{\textcolor[rgb]{0.64,0.00,0.00}{\textbf{#1}}}
\newcommand{\ExtensionTok}[1]{#1}
\newcommand{\FloatTok}[1]{\textcolor[rgb]{0.00,0.00,0.81}{#1}}
\newcommand{\FunctionTok}[1]{\textcolor[rgb]{0.13,0.29,0.53}{\textbf{#1}}}
\newcommand{\ImportTok}[1]{#1}
\newcommand{\InformationTok}[1]{\textcolor[rgb]{0.56,0.35,0.01}{\textbf{\textit{#1}}}}
\newcommand{\KeywordTok}[1]{\textcolor[rgb]{0.13,0.29,0.53}{\textbf{#1}}}
\newcommand{\NormalTok}[1]{#1}
\newcommand{\OperatorTok}[1]{\textcolor[rgb]{0.81,0.36,0.00}{\textbf{#1}}}
\newcommand{\OtherTok}[1]{\textcolor[rgb]{0.56,0.35,0.01}{#1}}
\newcommand{\PreprocessorTok}[1]{\textcolor[rgb]{0.56,0.35,0.01}{\textit{#1}}}
\newcommand{\RegionMarkerTok}[1]{#1}
\newcommand{\SpecialCharTok}[1]{\textcolor[rgb]{0.81,0.36,0.00}{\textbf{#1}}}
\newcommand{\SpecialStringTok}[1]{\textcolor[rgb]{0.31,0.60,0.02}{#1}}
\newcommand{\StringTok}[1]{\textcolor[rgb]{0.31,0.60,0.02}{#1}}
\newcommand{\VariableTok}[1]{\textcolor[rgb]{0.00,0.00,0.00}{#1}}
\newcommand{\VerbatimStringTok}[1]{\textcolor[rgb]{0.31,0.60,0.02}{#1}}
\newcommand{\WarningTok}[1]{\textcolor[rgb]{0.56,0.35,0.01}{\textbf{\textit{#1}}}}
\usepackage{graphicx}
\makeatletter
\newsavebox\pandoc@box
\newcommand*\pandocbounded[1]{% scales image to fit in text height/width
  \sbox\pandoc@box{#1}%
  \Gscale@div\@tempa{\textheight}{\dimexpr\ht\pandoc@box+\dp\pandoc@box\relax}%
  \Gscale@div\@tempb{\linewidth}{\wd\pandoc@box}%
  \ifdim\@tempb\p@<\@tempa\p@\let\@tempa\@tempb\fi% select the smaller of both
  \ifdim\@tempa\p@<\p@\scalebox{\@tempa}{\usebox\pandoc@box}%
  \else\usebox{\pandoc@box}%
  \fi%
}
% Set default figure placement to htbp
\def\fps@figure{htbp}
\makeatother
\setlength{\emergencystretch}{3em} % prevent overfull lines
\providecommand{\tightlist}{%
  \setlength{\itemsep}{0pt}\setlength{\parskip}{0pt}}
\usepackage{bookmark}
\IfFileExists{xurl.sty}{\usepackage{xurl}}{} % add URL line breaks if available
\urlstyle{same}
\hypersetup{
  pdftitle={Pooling Risk Ratios},
  hidelinks,
  pdfcreator={LaTeX via pandoc}}

\title{Pooling Risk Ratios}
\author{}
\date{\vspace{-2.5em}2026-04-17}

\begin{document}
\maketitle

\subsection{Introduction}\label{introduction}

Risk ratios (RR) --- also called relative risks --- are more
interpretable than odds ratios for common outcomes; they directly
express the ratio of event probabilities between two arms. RR
approximates OR only when events are rare (probability below
\textasciitilde10 \%); for common outcomes, the two diverge
substantially. Meta-analytic pooling of binary outcomes from cohort
studies and randomised trials therefore typically proceeds on the log-RR
scale, analogously to log-OR pooling.

\subsection{Prerequisites}\label{prerequisites}

A working understanding of risk ratios, log transformations for ratio
measures, and inverse-variance meta-analysis under random effects.

\subsection{Theory}\label{theory}

For a study with \(a_i\) events in \(a_i + b_i = n_{1i}\) exposed and
\(c_i\) events in \(c_i + d_i = n_{2i}\) unexposed:

\[\log \mathrm{RR}_i = \log \frac{a_i / n_{1i}}{c_i / n_{2i}}, \qquad \mathrm{Var}(\log \mathrm{RR}_i) = \frac{1}{a_i} - \frac{1}{n_{1i}} + \frac{1}{c_i} - \frac{1}{n_{2i}}.\]

Inverse-variance pooling on the log-RR scale follows the same
fixed-effect or random-effects distinction as any other meta-analysis.
Back-transformation to the natural RR scale via \(\exp\) produces
interpretable summaries with confidence intervals that respect the
multiplicative scale.

\subsection{Assumptions}\label{assumptions}

Studies are independent; log-RR is approximately Normal (typically a
good approximation when each study has at least 5 events per arm); the
population RR is the same across studies (fixed-effect) or distributed
around a mean with between-study variance \(\tau^2\) (random-effects).

\subsection{R Implementation}\label{r-implementation}

\begin{Shaded}
\begin{Highlighting}[]
\FunctionTok{library}\NormalTok{(metafor)}

\FunctionTok{set.seed}\NormalTok{(}\DecValTok{2026}\NormalTok{)}
\NormalTok{k }\OtherTok{\textless{}{-}} \DecValTok{10}
\NormalTok{df }\OtherTok{\textless{}{-}} \FunctionTok{data.frame}\NormalTok{(}
  \AttributeTok{ai =} \FunctionTok{rpois}\NormalTok{(k, }\DecValTok{30}\NormalTok{), }\AttributeTok{n1i =} \FunctionTok{rep}\NormalTok{(}\DecValTok{100}\NormalTok{, k),}
  \AttributeTok{ci =} \FunctionTok{rpois}\NormalTok{(k, }\DecValTok{20}\NormalTok{), }\AttributeTok{n2i =} \FunctionTok{rep}\NormalTok{(}\DecValTok{100}\NormalTok{, k)}
\NormalTok{)}

\NormalTok{es }\OtherTok{\textless{}{-}} \FunctionTok{escalc}\NormalTok{(}\AttributeTok{measure =} \StringTok{"RR"}\NormalTok{,}
             \AttributeTok{ai =}\NormalTok{ ai, }\AttributeTok{n1i =}\NormalTok{ n1i, }\AttributeTok{ci =}\NormalTok{ ci, }\AttributeTok{n2i =}\NormalTok{ n2i,}
             \AttributeTok{data =}\NormalTok{ df)}

\NormalTok{re }\OtherTok{\textless{}{-}} \FunctionTok{rma}\NormalTok{(yi, vi, }\AttributeTok{data =}\NormalTok{ es, }\AttributeTok{method =} \StringTok{"REML"}\NormalTok{)}
\FunctionTok{cat}\NormalTok{(}\StringTok{"Pooled RR:"}\NormalTok{, }\FunctionTok{round}\NormalTok{(}\FunctionTok{exp}\NormalTok{(re}\SpecialCharTok{$}\NormalTok{b), }\DecValTok{2}\NormalTok{),}
    \StringTok{" 95\% CI:"}\NormalTok{, }\FunctionTok{paste}\NormalTok{(}\FunctionTok{round}\NormalTok{(}\FunctionTok{exp}\NormalTok{(}\FunctionTok{c}\NormalTok{(re}\SpecialCharTok{$}\NormalTok{ci.lb, re}\SpecialCharTok{$}\NormalTok{ci.ub)), }\DecValTok{2}\NormalTok{), }\AttributeTok{collapse =} \StringTok{"{-}"}\NormalTok{), }\StringTok{"}\SpecialCharTok{\textbackslash{}n}\StringTok{"}\NormalTok{)}
\FunctionTok{cat}\NormalTok{(}\StringTok{"tau\^{}2:"}\NormalTok{, }\FunctionTok{round}\NormalTok{(re}\SpecialCharTok{$}\NormalTok{tau2, }\DecValTok{3}\NormalTok{), }\StringTok{"  I\^{}2:"}\NormalTok{, }\FunctionTok{round}\NormalTok{(re}\SpecialCharTok{$}\NormalTok{I2, }\DecValTok{1}\NormalTok{), }\StringTok{"\%}\SpecialCharTok{\textbackslash{}n}\StringTok{"}\NormalTok{)}
\end{Highlighting}
\end{Shaded}

\subsection{Output \& Results}\label{output-results}

\texttt{escalc(measure\ =\ "RR")} computes per-study log-RR and its
variance. \texttt{rma()} fits the random-effects model on the log scale;
back-transformation gives the pooled RR with confidence intervals on the
natural scale, plus \(\tau^2\) and \(I^2\) for heterogeneity.

\subsection{Interpretation}\label{interpretation}

A reporting sentence: ``Random-effects meta-analysis of 10 studies
(REML, \(n\) total = 2\{,\}000) gave a pooled risk ratio of 1.48 (95 \%
CI 1.24 to 1.77) with moderate heterogeneity (\(I^2 = 35 \%\),
\(\tau^2 = 0.02\)); the intervention is associated with a 48 \% higher
event rate, with the 95 \% CI excluding the null.'' Always report the
heterogeneity statistics alongside the pooled estimate.

\subsection{Practical Tips}\label{practical-tips}

\begin{itemize}
\tightlist
\item
  Use RR for cohort studies and randomised controlled trials, where
  baseline risk is meaningful; use OR for case-control studies, where RR
  is not estimable from the design.
\item
  For rare events (event rate below 10 \%), RR ≈ OR; for common events,
  the two diverge and the choice matters substantively.
\item
  Zero-event studies require a continuity correction
  (\texttt{escalc(...,\ add\ =\ 0.5)}) or a rare-events method (Peto OR,
  Mantel-Haenszel without correction); standard log-RR is undefined when
  any cell is zero.
\item
  Present forest plots and reports on the natural scale (RR), not the
  log scale; readers reason multiplicatively about ratios.
\item
  Mantel-Haenszel pooling (\texttt{metafor::rma.mh()}) is a common
  alternative for sparse-data settings; it does not require
  log-transformation and handles zero events more gracefully.
\item
  For individual-participant-data meta-analysis, fit a logistic
  regression with study as a random effect; this is more powerful than
  aggregate-data pooling when the IPD are available.
\end{itemize}

\subsection{R Packages Used}\label{r-packages-used}

\texttt{metafor::escalc()} with \texttt{measure\ =\ "RR"} and
\texttt{rma()} for the canonical workflow; \texttt{meta::metabin()} as
an alternative interface with cleaner forest-plot output;
\texttt{metafor::rma.mh()} for Mantel-Haenszel pooling;
\texttt{metafor::rma.peto()} for the Peto method on rare events.

\subsection{For Reviewers}\label{for-reviewers}

\textbf{What to look for in a paper using this method.}

\begin{itemize}
\tightlist
\item
  \textbf{Common misapplications.}
\item
  \textbf{Diagnostics that should be reported but often aren't.}
\item
  \textbf{Red flags in tables and figures.}
\item
  \textbf{What to verify.}
\item
  \textbf{What an adequate Methods paragraph must contain.}
\end{itemize}

\subsection{See also --- labs in this
chapter}\label{see-also-labs-in-this-chapter}

\begin{itemize}
\tightlist
\item
  \href{labs/lab-systematic-reviews-and-prisma.Rmd}{Systematic reviews
  and PRISMA}
\item
  \href{labs/lab-meta-analysis-basics.Rmd}{Meta-analysis basics}
\item
  \href{labs/lab-network-meta-analysis.Rmd}{Network meta-analysis}
\end{itemize}

\end{document}
